\section{Configuración experimental}\label{sec:experimentos}

\subsection{Modelo y hardware}

El orquestador de los tres agentes es \code{mistralai/Mistral-7B-Instruct-v0.3}~\cite{mistral}, cuantizado a 4 bits (NF4~\cite{qlora}, con cómputo en \code{float16}) mediante \code{bitsandbytes}.
Se ejecuta con \code{transformers} en una GPU Tesla T4 (16\,GB) de Google Colab.
La decodificación es voraz (\code{do\_sample=False}), con hasta 512 tokens nuevos, y se detiene en la cadena \code{Observation:}.
La corrida completa, de 100 registros (50 consultas con dos métodos), tomó unos 96 minutos.

\subsection{Sistemas comparados}

\begin{itemize}
  \item \textbf{Multi-agente.} Es el grafo de la sección~\ref{sec:arquitectura}, con \code{max\_rounds}~$=3$ y hasta 6 pasos ReAct por ronda.
  \item \textbf{Naive RAG.} Recupera 3 fragmentos de la misma base ChromaDB y hace una sola llamada al mismo modelo, sin herramientas ni auditoría.
  En GSM8K y MBPP la base de WikiQA no aporta contexto útil, así que ahí la comparación es contra un modelo que razona sin herramientas.
\end{itemize}

Ambos sistemas reciben exactamente las mismas 50 preguntas y los evalúa el mismo juez.

\subsection{Métricas}

Sea $q$ una pregunta, $a$ la respuesta y $C$ el contexto del método: los fragmentos recuperados y, en el multi-agente, también las salidas de \code{calc} y \code{python}.

\paragraph{Faithfulness.}
El juez descompone $a$ en a lo sumo 8 afirmaciones atómicas $s_1,\dots,s_n$ y marca cada una como inferible (\code{SI}) o no (\code{NO}) a partir de $C$ y de los datos de $q$:
\begin{equation}
  \text{Faithfulness}(a, C) = \frac{\left|\{\,s_i : s_i \text{ se infiere de } C \cup q\,\}\right|}{n}.
\end{equation}
Los datos del enunciado cuentan como información dada porque en GSM8K y MBPP la respuesta los repite con razón.

\paragraph{Answer Relevance.}
El juez asigna un entero $r \in \{0,\dots,10\}$ según qué tan directamente $a$ responde $q$, sin juzgar si es verdadera, y se normaliza como $r/10$.
Es una simplificación de RAGAS, que genera preguntas a partir de la respuesta y mide su similitud con la original.

\paragraph{Verificación objetiva.}
En GSM8K se extrae el número que sigue a \code{Respuesta final:} (o, si no existe esa línea, el último número del texto) y se compara con la referencia.
En MBPP se ejecuta el código de la respuesta con sus \code{assert}.
Sobre estas 30 consultas se reportan tres métricas:
\begin{align}
  \text{exactitud} &= \frac{\#\,\text{correctas}}{\#\,\text{consultas}}, \\
  \text{exactitud al responder} &= \frac{\#\,\text{correctas}}{\#\,\text{respondidas}}, \\
  \text{incorrectas entregadas} &= \frac{\#\,\text{respondidas e incorrectas}}{\#\,\text{consultas}}.
\end{align}
La última mide cuántas alucinaciones llegan al usuario: una abstención no cuenta como incorrecta.

\paragraph{Abstenciones.}
Si el Auditor nunca aprueba, el multi-agente no responde.
Una abstención puntúa 0 en \emph{Faithfulness} y \emph{Answer Relevance}, y cuenta como no correcta en la exactitud.
Para separar ``qué tan fiel es cuando responde'' de ``cuántas veces se abstiene'', se reporta además la fidelidad promediada solo sobre las respuestas entregadas.

\paragraph{Juez.}
El juez es el mismo Mistral-7B que orquesta.
Así la evaluación cabe en la misma GPU, pero el juez es ruidoso y podría favorecer sus propias respuestas.
Si el juez no devuelve el formato pedido se reintenta una vez; si vuelve a fallar, la métrica queda vacía y se cuenta aparte.

\paragraph{Prueba estadística.}
Para la exactitud se usa la prueba exacta de McNemar sobre los pares discordantes: las consultas que acierta un solo método.
Es la prueba adecuada para dos clasificadores evaluados sobre las mismas muestras.
